\documentclass[10pt,a4paper]{amsart}
\usepackage[margin=19mm]{geometry}
\usepackage{amsmath,amssymb,mathtools}
\usepackage[scr=rsfs,cal=euler]{mathalfa}
\usepackage{xfrac}
\usepackage[unicode,hidelinks,bookmarks=false]{hyperref}
\newcommand{\ie}{i.e.\ }
\newcommand{\cf}{cf.\ }
\newcommand{\resp}{respectively\ }
\newcommand{\Subsection}{\subsection}
\providecommand{\nobreakdash}{\nobreak\hbox{-}}
\setlength{\emergencystretch}{2em}
\multlinegap0pt
\usepackage{amssymb}
\usepackage{tikz}
\usepackage{bm}
\usepackage[shortlabels]{enumitem}
\newtheorem{theorem}{Theorem}
\newcommand{\Fib}{{\mathcal F}}
\title{Munarini graphs: a generalization of Fibonacci cubes and Pell graphs. Part I}
\author{Michel Mollard}
\date{Source: arXiv:2605.14613v1 (CC BY 4.0)}
\begin{document}
\maketitle

\noindent\emph{Source and licence.} Munarini graphs: a generalization of Fibonacci cubes and Pell graphs. Part I. Version arXiv:2605.14613v1 (CC BY 4.0), distributed by arXiv under the Creative Commons Attribution 4.0 International licence, http://creativecommons.org/licenses/by/4.0/, as recorded in the arXiv licence metadata for this exact version and shown on the abstract page https://arxiv.org/abs/2605.14613v1. Full licence notice: \texttt{licenses/arxiv-2605.14613-CC-BY-4.0.txt}.\par\smallskip

\noindent\textit{Editorial scope.} Counted unit: the article's theorem th:isoembed (source0.tex lines 886--888). The article's complete original proof of it is delivered with it (source0.tex lines 889--906, one \texttt{proof} environment of 18 lines). No mathematics is written, repaired or re-derived: the statement and the proof are the article's own lines, in the article's own notation, and no retained line is substituted or re-flowed.  No further source-local context is needed: every cross-reference of every retained line resolves inside the two delivered spans.  Nothing was omitted from any retained span, so the package carries no omission record.  No retained line contains a citation, so the wrapper carries no bibliography and no bibliography entry.  No retained line contains a figure environment, an image-inclusion command or drawing code, so no image is delivered and the package declares no asset dependency.  Editorial formatting only, in addition: an AMS \texttt{amsart} preamble that restates the article's own declarations and macros used by the retained lines, namely the environments \texttt{theorem} on separate counters, together with the standard packages those lines need.  The retained environments print in this document as Theorem 1 in order of appearance, whereas the article prints the same items with the numbering of its own full document.  The complete native source of the article as distributed in the arXiv e-print for this exact version is bundled unchanged as \texttt{sources/200599/source0.tex}.
\begin{theorem}\label{th:isoembed}
$M_{n,k}$ is isomorphic to $\Gamma_{n,k}^\star$, and in particular $M_{n,k}$ is a subgraph of $\Gamma_{kn-1}$.
\end{theorem}

\begin{proof}
Let $\Psi(i)=A_i$ for $i\in\{0,\ldots,k-1\}$ and $\Psi(kk)=A_kA_0$. The mapping $\Psi$ between the alphabets $\{0,1,\ldots,k-1,kk\}$ and $\{A_0,A_1,\ldots,A_{k-1},A_kA_0\}$  extends naturally to a bijection $\widehat{\Psi}$ between the monoids they generate. Since $\widehat{\Psi}$ maps a $(k+1)$-ary string of length $n$ to a binary string of length $kn$, $\widehat{\Psi}$ is therefore a bijection between $\Fib_{n,k}$ and $\Fib_{n,k}^\star$. 

Let $i,j\in\{0,1,\ldots,k\}$. Then, there exists an edge in $\Gamma_k$ between $A_i$ and $A_j$ if and only if $i=0$ and $j\in\{1,\ldots,k\}$ (or vice versa). 

Assume there exists an edge in $M_{n,k}$ between $u$ and $v$. Then two cases arise.

In the first case,  $u$ and $v$ can be written as  $u=a0b$ and $v=aib$ with $i\in\{1,\ldots,k-1\}$ where $a$ and $b$ are generalized Pell strings. We then have an edge in $\Gamma_{n,k}^\star$ between 
$\widehat{\Psi}(u)=\widehat{\Psi}(a)A_0\widehat{\Psi}(b)$ and $\widehat{\Psi}(v)=\widehat{\Psi}(a)A_i\widehat{\Psi}(b)$.

The second possibility is that $u=a00b$ and $v=akkb$. Then there exists an edge in $\Gamma_{n,k}^\star$ between 
$\widehat{\Psi}(u)=\widehat{\Psi}(a)A_kA_0\widehat{\Psi}(b)$ and $\widehat{\Psi}(v)=\widehat{\Psi}(a)A_0A_0\widehat{\Psi}(b)$. Note that no other adjacency is possible in $\Gamma_{n,k}^\star$ in this case, because for $i\neq0$, $\widehat{\Psi}(a)A_kA_i\widehat{\Psi}(b)$ does not belong to $\Fib_{n,k}^\star$. 
  
Conversely, assume there exists an edge in $\Gamma_{n,k}^\star$	between $xA_iy$ and $xA_0y$, $i\in\{1,\ldots,k-1\}$ , where $x,y$ are  Munarini strings. Then there exists an edge in $M_{n,k}$ between $\widehat{\Psi}^{-1}(xA_iy)=\widehat{\Psi}^{-1}(x)i\widehat{\Psi}^{-1}(y)$ and $\widehat{\Psi}^{-1}(xA_0y)=\widehat{\Psi}^{-1}(x)0\widehat{\Psi}^{-1}(y)$. The case of an edge between $xA_kA_0y$ and $xA_0A_0y$ is similar since $\widehat{\Psi}^{-1}(A_kA_0)=kk$ and $\widehat{\Psi}^{-1}(A_0A_0)=00$.
$\widehat{\Psi}$ is therefore a graph isomorphism between $M_{n,k}$ and $\Gamma_{n,k}^\star$ .

Since all strings in $\Fib_{n,k}^\star$ end with 0, removing this 0 makes $\Gamma_{n,k}^\star$ itself isomorphic to a subgraph of $\Gamma_{kn-1}$.
\end{proof}

\end{document}
